% ============================================================================
\section{Experimental design}\label{sec:design}
% ============================================================================

\subsection{Network, products and costs}

The base network consists of a depot and the first ten customers of Solomon's R101 instance \citep{solomon1987algorithms}, which are scattered around the depot at distances of 11 to 32; coordinates define Euclidean distances and time windows are ignored. Two further layouts use the first ten customers of C101, a tight cluster 15--21 distance units from the depot (\emph{clustered}), and of RC101, a remote cluster of eight stores 33--40 units away plus two stores near the depot (\emph{mixed}). Two vehicles serve the stores. Vehicle capacity is set through the fleet ratio $\rho=KQ/\bar D$, where $\bar D$ is mean total daily demand; the base case has $\rho=1.25$ (\VehicleCap{} units per vehicle), so that the fleet can carry a quarter more than average demand, as is common in grocery distribution where vehicles are sized for peak days. Shelf capacity is three days of mean store demand. The base product has a shelf life of $m=3$ days. Costs are $h=1$ per unit and night, $p_s=5$ per expired unit and $p_u=20$ per lost sale, with routing cost one per distance unit; per unit of demand, waste is five times and a lost sale twenty times as expensive as one night of holding.

\subsection{Demand and information}

Demand in the synthetic settings is generated by a contextual negative-binomial model (code in the replication package). Store base levels range from 10 to 22 units per day; demand varies with the day of the week (Saturday $+30\%$), announced promotions (12\% of store-days, discounts of 10--30\%, multiplying demand by up to 1.8 and lowering the negative-binomial size parameter from 8 to 3.2, so that promotions raise both level and variability), temperature, holidays, and a latent store level following a log-AR(1) process with autocorrelation 0.92 that can only be inferred from recent sales. The coefficient of variation of daily store demand is about 0.4 on regular days. The temperature forecast error grows with the lead time. The model is designed to reproduce qualitative features of fresh retail demand (weekly pattern, promotion lifts, heteroscedasticity, persistence), not calibrated to a specific retailer; the real-data case below tests the policies on observed demand.

Each experiment uses 200 training, 50 validation and 100 test episodes, generated with disjoint random seeds. An episode consists of 35 days of history (to compute lagged features), 21 decision days, and four further days so that lookahead forecasts and the perfect-information reference are defined on the last decision days. Every policy starts each test episode with empty stores and is evaluated on the same 100 demand paths (common random numbers).

\subsection{Real-data case}\label{sec:realdata}

The real-data case replaces the synthetic demand by observed sales of FreshRetailNet-50K \citep{wang2025freshretailnet}, a public dataset of 50,000 store--product series of a fresh-food retailer over 97 days (its 90-day training and seven-day evaluation parts combined), with hourly sales, hourly out-of-stock indicators, discounts, holidays and weather (CC BY 4.0 licence). An episode is one product sold at ten stores of the same city over the last 60 days of the data (35 history, 21 decision and four further days). We keep series without missing days and with fewer than 10\% zero-sales days and form at most one ten-store group per city and product. Products, not groups, are assigned to the training, validation and test sets, so that no product appears in two of them: the 200 training, 50 validation and 100 test episodes come from 101, 28 and 56 products in 13 cities.

Three preparation steps are needed and disclosed. First, sales are censored on days with stock-outs, which concern about 40\% of the store-days. Latent demand is estimated by dividing a day's sales by the share of the series' typical intraday sales profile that falls into in-stock hours, the profile being estimated on the series' fully in-stock days and the share floored at 0.25. Second, the products differ in volume by orders of magnitude, whereas the network, fleet and shelf capacities are those of the synthetic study; each store's latent demand is therefore rescaled to the store's base volume using the mean over the 35 history days only, so that no information from the decision days is used. Third, fleet and shelf capacities are dimensioned on the network's average volume, so groups whose sales level changes by more than $-20\%/{+25\%}$ between the history and decision windows (typically products entering or leaving their season) are excluded; in them the fleet, not the replenishment policy, would determine the outcome. This selection concerns only which data enter the experiment; forecasts and policies use past information only. In the resulting test episodes demand on the decision days is on average 6\% above the history level, the coefficient of variation of daily store demand is 0.55, and discounts are offered on 44\% of the store-days. Forecast features are those of the synthetic study, with the observed temperature used as its own forecast.

\subsection{Configurations}

Besides the base case we vary one factor at a time: network layout (clustered, mixed); shelf life $m\in\{2,4\}$; lost-sales penalty $p_u\in\{10,40\}$; fleet ratio $\rho\in\{1.10,1.60\}$; routing cost $c^{\mathrm{r}}=3$ per distance unit, which makes consolidation of deliveries more valuable; demand variability (negative-binomial size 3 instead of 8, i.e.\ coefficient of variation about 0.6 on regular days); partial LIFO issuing, where 30\% of customers take the freshest units; and the real-data case. In the LIFO configuration all policies are tuned and trained under FIFO and evaluated under mixed issuing, which tests robustness to a misspecified issuing assumption that practitioners often face.

\subsection{Tuning, training and inference}

All parameters are selected on validation episodes by closed-loop simulation: BSP and BSP-EW over $R\in\{1,2\}$ and $\alpha\in\{0.6,0.7,0.8,0.9,0.95\}$; UL$_\lambda$ over $\lambda\in\{1,2,3\}$, $\alpha\in\{0.7,0.8,0.9,0.95\}$, $\beta\in\{0.5,0.7,0.9\}$; LA-Quantile over $\alpha\in\{0.5,\dots,0.9\}$; LA-Quantile2 over $\alpha_0\in\{0.9,0.95\}$ and $\alpha_1\in\{0.6,0.7,0.8\}$; LA-Robust over $\Gamma\in\{0.25,0.5,1,1.5\}$.

UL-DFL is trained on all 200 training episodes for 30 epochs with Adam (learning rate $10^{-3}$, minibatches of ten episodes, gradient norm clipped at 5, $L_2$ penalty $10^{-4}$ on the output layer), from three random seeds in every configuration; NN-DFL is trained identically but for 60 epochs. LA-DFL is trained on the decision states visited by the tuned LA-Quantile policy in 100 training episodes (2,100 states), with eight epochs, visit patterns re-optimised every second epoch, Adam (learning rate $2\cdot10^{-3}$, minibatch 64), a two-layer network with 64 SiLU units, and three seeds in the base case and one elsewhere. For all learners the network and the seed with the lowest closed-loop validation cost are kept. NN-DFL was added after the primary hypotheses had been fixed; its comparisons are reported as secondary evidence. Test episodes are used once.

We prespecify two families of primary hypotheses in the base case. For UL-DFL they compare it with UL$_\lambda$, BSP-EW, ScenPO-SS, LA-DFL and LA-Quantile2; for LA-DFL with LA-Quantile, ScenPO-SS, LA-Robust, BSP-EW and UL$_\lambda$. The unit of analysis is the episode; for each comparison we report the mean paired cost difference, a 95\% paired bootstrap interval (10,000 resamples), paired $t$ and Wilcoxon signed-rank tests, and Holm-adjusted $p$-values within each family \citep{holm1979simple}. All other comparisons and configurations are reported with paired intervals as secondary evidence.

\subsection{Optimisation settings}\label{sec:solver}

Lookahead MILPs are solved by HiGHS \citep{huangfu2018parallelizing} through SciPy with a relative gap of 0.5\% and time limits of 20 seconds (deterministic) and 60 seconds (scenario), with integer vehicle assignments on the implemented day and relaxed assignments on later lookahead days (Section~\ref{sec:lookahead}). The LA-DFL oracle uses the same model with a 1\% gap. If a time limit were reached the best incumbent would be used and its proven gap recorded; Section~\ref{sec:computation} reports solver statistics. All experiments run on a 16-core laptop computer (Intel Core Ultra 9 285H) with 13 to 15 parallel worker processes; the code, seeds and configuration files are in the replication package.
